\scanpage{13}
que $A'$ est muni d'une \uline{donnée de descente} relativement à $S'\unsure{-|\to}S$. En particulier, \deleted{\unsure{l|E}} $ {}_{\ell}A'$ est muni d'une telle donnée de descente,. \deleted{qui} \deleted{est} Or $S'-S$ est \inserted{morphisme} de descente effective pour les schémas relatifs finis étales (SGA 1 VIII), donc on trouve un \deleted{\unsure{se|tel}} tel \unsure{rev.|sch.|schéma} \inserted{en groupes} \deleted{\unsure{B|R}} $N$ de $S$, tel que $N\mid U\simeq{}_{\ell}(A_U)$. Cela montre que l'on peut trouver un morphisme \inserted{\uline{fini}} \uline{étale} surjectif $S_1\to S$ qui splitte $ {}_{\ell}(A_U)$, \deleted{Remplaçant alors $S'$} et on peut supposer $S'$ \deleted{fini} \uline{étale} sur $S$, donc plat. Alors ($A'$ étant projectif sur $S'$) on \deleted{trouve une donnée de \unsure{des}} peut \uline{descendre} $A'$ en $A$, qui résoud le pb.

d) Cas $S$ réduit, \inserted{\deleted{la \unsure{normalisée}} tel qu'il existe $S'-S$ fini surjectif, avec $S'$ géom. unib.} $S'$ fini sur $S$. Utilisant le fait que $S$ est géom unibranche en les points de $S-U$, on voit que l'image inverse $U''=U'\times_U U'$ de $U$ dans $S''=S'\times_S S'$ est dense dans $S''$: \deleted{Par le cas \unsure{b|6}), on}

\diagram{\begin{tikzcd}[ampersand replacement=\&,column sep=small,row sep=small]
U \arrow[d,arrows={Hooks[right]-}] \& U' \arrow[l] \arrow[d,arrows={Hooks[right]-}] \& U'' \arrow[l,shift left=.45ex] \arrow[l,shift right=.45ex] \arrow[d,arrows={Hooks[right]-}]\\
S \& S' \arrow[l] \& S'' \arrow[l,dash,shift left=.45ex] \arrow[l,dash,shift right=.45ex]
\end{tikzcd}}

Par le cas c), on trouve sur $S'$ un schéma abélien $A'$, et par la pleine fidélité (pour $S'',U''$) une donnée de descente sur $A'$, relativement à $S'\to S$. Notons que le morphisme $S'-S$, étant un épimorphisme fini, est un morphisme de descente pour les schémas abéliens, \deleted{\unsure{\illegible{1}|E}} et aussi pour les gpes $p$-divisibles. Donc on est ramené à
% slanted
\margin{lemme à dégager.}
% archive sheet number: 12
